\scanpage{21}
\[
\xi'''\in\operatorname{Hom}(N,J)/\operatorname{Im}\operatorname{Hom}(M_{v_0},J)
\]

\deleted{Cet} Cet élément localisé en $\eta$ est nul, \deleted{\unsure{car}} et c'est là tout ce que nous pouvons tirer du fait que $u_0\circ\varphi_0$ et $u_{0\eta}$ se prolongent. Mais \uline{nous savons plus} : les prolongements \inserted{$v$ et $u_\eta$} peuvent se faire de façon que $v_\eta=\deleted{\unsure{\varphi_\eta}}u_\eta\varphi_\eta$. À priori, \deleted{\unsure{l'obs|l'idée}} une fois acquis le fait de la possibilité \uline{séparée} des deux obstructions, i.e. dans $(1)$ de l'existence de deux \unsure{élémts|éléments} dans les coins, la \deleted{\unsure{obs}} possibilité de les choisir de façon cohérente se trouve dans

\[
\begin{aligned}
&\operatorname{Hom}(M_{v_0},J_\eta)/\operatorname{Im}\big(\operatorname{Hom}(M_{v_0},J)+\operatorname{Hom}(M_{u_0},J_\eta)\big)\\
&\subset\operatorname{Hom}(N,J_\eta)/\operatorname{Im}\operatorname{Hom}(M_{v_0},J)
\end{aligned}
\]

Soit $\unsure{\gamma|y|\eta}$ cet élément. Notons que $\xi'''$ définit également un élément du dit quotient, puisque

\[
\operatorname{Hom}(N,J_{\deleted{\unsure{\eta|n}}})/\operatorname{Im}\operatorname{Hom}(M_{v_0},J)\subset\operatorname{Hom}(N,J_\eta)/\operatorname{Im}\ldots
\]

Nous admettrons que l'on a $\gamma=\xi'''$. Comme par hypothèse $\unsure{\gamma|\eta}'=0$, on a $\xi'''=0$, OK.

% archive sheet number: 20
